% GENERATED FILE. Edit canonical lessons, then run npm run generate:books.
% canonical-source-hash: fnv1a64:804f1694066e38c1
% canonical-lessons: ML-C514-varalcca, ML-C514-vintukiruka, ML-C514-kariyuka, ML-C514-murinnupokuka, ML-C514-nilaykkuka

\chapter{Drought and Power Cuts}
\label{ch:ml-drought-and-power-cuts}
\hlchaptermodality{514}

\begin{chapteropening}
\textbf{By the end of this chapter:} \emph{I can say a drought, to crack open, to get scorched, to be disconnected and to stop working.}

Complete the last lesson of chapter 514: I can say a drought, to crack open, to get scorched, to be disconnected and to stop working.
\end{chapteropening}

\section[\texorpdfstring{\ml{വരൾച്ച} (varaḷcca)}{വരൾച്ച (varaḷcca)}]{\ml{വരൾച്ച} (varaḷcca) — a drought}

\label{lesson:ML-C514-varalcca}



\begin{quote}
Before the new one: say the Malayalam for to be uprooted, then the Malayalam for to be in distress.
\end{quote}

\subsection*{You'll want to know: \ml{വരൾച്ച}}
\textbf{\ml{വരൾച്ച}} — \emph{varaḷcca} — \textquotedblleft{}a drought\textquotedblright{}.

\begin{quote}
\ml{വരൾച്ച} \ml{കാരണം} \ml{കർഷകർ} \ml{ബുദ്ധിമുട്ടി}.
\end{quote}

\begin{quote}
\emph{varaḷcca kāraṇaṁ karṣakar buddhimuṭṭi.}
\end{quote}

\begin{quote}
The farmers suffered because of the drought.
\end{quote}

The noun of \ml{വരളുക}, to dry up.

\begin{cousinweb}
\textbf{\ml{വരൾച്ച}} (\emph{varaḷcca}) is native Dravidian, the noun of \ml{വരളുക}, to grow dry (DEDR 5320).
\end{cousinweb}

\begin{grammarlens}[title={What you've built}]
One more word for getting around.
\end{grammarlens}

\subsection*{Your turn}
\begin{itemize}
\raggedright
  \item \emph{Say it:} \emph{varaḷcca}
  \item \emph{Say it:} \emph{varaḷcca}, once more
  \item \emph{Say it:} \emph{varaḷcca kāraṇaṁ karṣakar buddhimuṭṭi.} — The farmers suffered because of the drought.
  \item \emph{Recall:} say the Malayalam for to be uprooted, then the Malayalam for to be in distress, then say \emph{varaḷcca} again
\end{itemize}

\begin{tcolorbox}[breakable,colback=teal!4,colframe=teal!35!black,title={Before you move on}]
What does \ml{വരൾച്ച} mean? (\textquotedblleft{}A drought\textquotedblright{}.) Say it once more.
\end{tcolorbox}
\section[\texorpdfstring{\ml{വിണ്ടുകീറുക} (viṇṭukīṟuka)}{വിണ്ടുകീറുക (viṇṭukīṟuka)}]{\ml{വിണ്ടുകീറുക} (viṇṭukīṟuka) — to crack open}

\label{lesson:ML-C514-vintukiruka}



\begin{quote}
Before the new one: say the Malayalam for to be in distress, then the Malayalam for a drought.
\end{quote}

\subsection*{You'll want to know: \ml{വിണ്ടുകീറുക}}
\textbf{\ml{വിണ്ടുകീറുക}} — \emph{viṇṭukīṟuka} — \textquotedblleft{}to crack open\textquotedblright{}.

\begin{quote}
\ml{വരൾച്ചയിൽ} \ml{വയലുകൾ} \ml{വിണ്ടുകീറി}.
\end{quote}

\begin{quote}
\emph{varaḷccayil vayalukaḷ viṇṭukīṟi.}
\end{quote}

\begin{quote}
In the drought the fields cracked open.
\end{quote}

Past: \ml{വിണ്ടുകീറി}, \emph{viṇṭukīṟi}.

\ml{വിണ്ട്}, cracking, plus \ml{കീറുക}, to tear.

\begin{cousinweb}
\textbf{\ml{വിണ്ടുകീറുക}} (\emph{viṇṭukīṟuka}) is a native Dravidian compound built on \ml{വിള്ളുക}, to burst open or crack (DEDR 5432), and \ml{കീറുക}, to tear.
\end{cousinweb}

\begin{grammarlens}[title={What you've built}]
One more word for getting around.
\end{grammarlens}

\subsection*{Your turn}
\begin{itemize}
\raggedright
  \item \emph{Say it:} \emph{viṇṭukīṟuka}
  \item \emph{Say it:} \emph{viṇṭukīṟuka}, once more
  \item \emph{Say it:} \emph{varaḷccayil vayalukaḷ viṇṭukīṟi.} — In the drought the fields cracked open.
  \item \emph{Recall:} say the Malayalam for to be in distress, then the Malayalam for a drought, then say \emph{viṇṭukīṟuka} again
\end{itemize}

\begin{tcolorbox}[breakable,colback=teal!4,colframe=teal!35!black,title={Before you move on}]
What does \ml{വിണ്ടുകീറുക} mean? (\textquotedblleft{}To crack open\textquotedblright{}.) Say it once more.
\end{tcolorbox}
\section[\texorpdfstring{\ml{കരിയുക} (kariyuka)}{കരിയുക (kariyuka)}]{\ml{കരിയുക} (kariyuka) — to get scorched; to wither}

\label{lesson:ML-C514-kariyuka}



\begin{quote}
Before the new one: say the Malayalam for a drought, then the Malayalam for to crack open.
\end{quote}

\subsection*{You'll want to know: \ml{കരിയുക}}
\textbf{\ml{കരിയുക}} — \emph{kariyuka} — \textquotedblleft{}to get scorched; to wither\textquotedblright{}.

\begin{quote}
\ml{കടുത്ത} \ml{വെയിലിൽ} \ml{ചെടികൾ} \ml{കരിഞ്ഞു}.
\end{quote}

\begin{quote}
\emph{kaṭutta veyilil ceṭikaḷ kariññu.}
\end{quote}

\begin{quote}
In the fierce sun the plants were scorched.
\end{quote}

Past: \ml{കരിഞ്ഞു}, \emph{kariññu}.

For plants and food: \ml{ചോറ്} \ml{കരിഞ്ഞു}, the rice got burnt.

\begin{cousinweb}
\textbf{\ml{കരിയുക}} (\emph{kariyuka}) is native Dravidian, \textquotedblleft{}to be scorched, to be singed\textquotedblright{}, from the same root as \ml{കരി}, charcoal (DEDR 1278).
\end{cousinweb}

\begin{grammarlens}[title={What you've built}]
One more word for getting around.
\end{grammarlens}

\subsection*{Your turn}
\begin{itemize}
\raggedright
  \item \emph{Say it:} \emph{kariyuka}
  \item \emph{Say it:} \emph{kariyuka}, once more
  \item \emph{Say it:} \emph{kaṭutta veyilil ceṭikaḷ kariññu.} — In the fierce sun the plants were scorched.
  \item \emph{Recall:} say the Malayalam for a drought, then the Malayalam for to crack open, then say \emph{kariyuka} again
\end{itemize}

\begin{tcolorbox}[breakable,colback=teal!4,colframe=teal!35!black,title={Before you move on}]
What does \ml{കരിയുക} mean? (\textquotedblleft{}To get scorched; to wither\textquotedblright{}.) Say it once more.
\end{tcolorbox}
\section[\texorpdfstring{\ml{മുറിഞ്ഞുപോകുക} (muṟiññupōkuka)}{മുറിഞ്ഞുപോകുക (muṟiññupōkuka)}]{\ml{മുറിഞ്ഞുപോകുക} (muṟiññupōkuka) — to be disconnected; to break off}

\label{lesson:ML-C514-murinnupokuka}



\begin{quote}
Before the new one: say the Malayalam for to crack open, then the Malayalam for to get scorched.
\end{quote}

\subsection*{You'll want to know: \ml{മുറിഞ്ഞുപോകുക}}
\textbf{\ml{മുറിഞ്ഞുപോകുക}} — \emph{muṟiññupōkuka} — \textquotedblleft{}to be disconnected; to break off\textquotedblright{}.

\begin{quote}
\ml{സംസാരിക്കുന്നതിനിടെ} \ml{ഫോൺ} \ml{മുറിഞ്ഞുപോയി}.
\end{quote}

\begin{quote}
\emph{saṁsārikkunnatiniṭe phōṇ muṟiññupōyi.}
\end{quote}

\begin{quote}
In the middle of talking, the phone call was cut off.
\end{quote}

Past: \ml{മുറിഞ്ഞുപോയി}, \emph{muṟiññupōyi}.

\ml{മുറിഞ്ഞ്}, being cut, plus \ml{പോകുക}: for calls, power and roads.

\begin{cousinweb}
\textbf{\ml{മുറിഞ്ഞുപോകുക}} (\emph{muṟiññupōkuka}) is a native Dravidian compound of \ml{മുറിയുക}, to be cut or broken (DEDR 5008), and \ml{പോകുക}, to go.
\end{cousinweb}

\begin{grammarlens}[title={What you've built}]
One more word for getting around.
\end{grammarlens}

\subsection*{Your turn}
\begin{itemize}
\raggedright
  \item \emph{Say it:} \emph{muṟiññupōkuka}
  \item \emph{Say it:} \emph{muṟiññupōkuka}, once more
  \item \emph{Say it:} \emph{saṁsārikkunnatiniṭe phōṇ muṟiññupōyi.} — In the middle of talking, the phone call was cut off.
  \item \emph{Recall:} say the Malayalam for to crack open, then the Malayalam for to get scorched, then say \emph{muṟiññupōkuka} again
\end{itemize}

\begin{tcolorbox}[breakable,colback=teal!4,colframe=teal!35!black,title={Before you move on}]
What does \ml{മുറിഞ്ഞുപോകുക} mean? (\textquotedblleft{}To be disconnected; to break off\textquotedblright{}.) Say it once more.
\end{tcolorbox}
\section[\texorpdfstring{\ml{നിലയ്ക്കുക} (nilaykkuka)}{നിലയ്ക്കുക (nilaykkuka)}]{\ml{നിലയ്ക്കുക} (nilaykkuka) — to stop working; to cease}

\label{lesson:ML-C514-nilaykkuka}



\begin{quote}
Before the new one: say the Malayalam for to get scorched, then the Malayalam for to be disconnected.
\end{quote}

\subsection*{You'll want to know: \ml{നിലയ്ക്കുക}}
\textbf{\ml{നിലയ്ക്കുക}} — \emph{nilaykkuka} — \textquotedblleft{}to stop working; to cease\textquotedblright{}.

\begin{quote}
\ml{വൈദ്യുതി} \ml{പോയപ്പോൾ} \ml{ഫാൻ} \ml{നിലച്ചു}.
\end{quote}

\begin{quote}
\emph{vaidyuti pōyappōḷ phān nilaccu.}
\end{quote}

\begin{quote}
When the power went, the fan stopped.
\end{quote}

Past: \ml{നിലച്ചു}, \emph{nilaccu}.

For machines, sounds and breath: \ml{ശബ്ദം} \ml{നിലച്ചു}, the sound stopped.

\begin{cousinweb}
\textbf{\ml{നിലയ്ക്കുക}} (\emph{nilaykkuka}) is native Dravidian, from the root of \ml{നിൽക്കുക}, to stand; DEDR gives it as to come to a stand, to cease (DEDR 3675).
\end{cousinweb}

\begin{grammarlens}[title={What you've built}]
One more word for getting around.
\end{grammarlens}

\subsection*{Your turn}
\begin{itemize}
\raggedright
  \item \emph{Say it:} \emph{nilaykkuka}
  \item \emph{Say it:} \emph{nilaykkuka}, once more
  \item \emph{Say it:} \emph{vaidyuti pōyappōḷ phān nilaccu.} — When the power went, the fan stopped.
  \item \emph{Recall:} say the Malayalam for to get scorched, then the Malayalam for to be disconnected, then say \emph{nilaykkuka} again
\end{itemize}

\begin{tcolorbox}[breakable,colback=teal!4,colframe=teal!35!black,title={Before you move on}]
What does \ml{നിലയ്ക്കുക} mean? (\textquotedblleft{}To stop working; to cease\textquotedblright{}.) Say it once more.
\end{tcolorbox}
